\section{Equilibrium multiplicity and a cost-floor benchmark}\label{sec:equilibrium}

The failure of the published profile does not by itself select a replacement equilibrium. The reason is that a foreclosed outsider can post a price that attracts no demand, and such a quote can alter the active firms' deviation incentives even though the outsider sells nothing. We therefore distinguish the price game stated in the published model from a separate price game that explicitly rules out below-cost prices.

Write a symmetric foreclosed pure-strategy profile in one union-member market as
\[
 (p_1,p_2,p_3)=(s,s,r),
\]
where the outsider has zero demand up to a zero-measure tie. The characterization below is deliberately limited to this class. It does not characterize asymmetric pure equilibria, mixed equilibria, or the government-stage equilibrium under continuation multiplicity.

At the outsider's ideal location, either member product is delivered at price $s+1$. Hence foreclosure requires $r\geq s+1$. In the unrestricted game, zero outsider demand can be a best response only if $c\geq s+1$: when $c<s+1$, the outsider can quote a price strictly between $c$ and $s+1$, obtain positive demand near its ideal location, and earn strictly positive profit.

\begin{proposition}[Pricing and multiplicity]\label{prop:equilibrium}
Consider the published quadratic-transport price game in the strict post-foreclosure range $5/2<c<5$.

\begin{enumerate}
\item The Appendix-B profile $(3/2,3/2,c)$ is not a Nash equilibrium.

\item Within symmetric, foreclosed, pure-strategy profiles $(s,s,r)$, the unrestricted original price game has a Nash equilibrium if and only if either
\begin{align*}
 &\frac32\leq s<2,\qquad r=s+1,\qquad c\geq s+1, \tag{U1}\\
 &s=2,\qquad r\geq3,\qquad c\geq3. \tag{U2}
\end{align*}
Thus the admissible member price spans
\begin{equation}\label{eq:unrestricted-member-set}
 \frac32\leq s\leq \min\{2,c-1\},
\end{equation}
subject to the outsider-price conditions in U1--U2. The unrestricted game therefore has genuine member-price multiplicity.

\item Define a separate cost-floor price game by imposing $p_i\geq mc_i$ in the market under study. Within symmetric, foreclosed, pure-strategy equilibria of that restricted game,
\begin{equation}\label{eq:piecewise-price}
 p_M(c)=
 \begin{cases}
 c-1, & \dfrac52<c<3,\\[4pt]
 2, & 3\leq c<5.
 \end{cases}
\end{equation}
For $5/2<c<3$, the complete symmetric foreclosed profile is $(c-1,c-1,c)$. For $3\leq c<5$, the member price is $2$ and any outsider quote $r\geq c$ is compatible with the stated equilibrium class.
\end{enumerate}
\end{proposition}

\begin{proof}
Part 1 follows from \eqref{eq:published-profile-gain}. We prove part 2 by necessity and sufficiency within the stated symmetric foreclosed class.

For U1, fix $p_2=s$ and $p_3=s+1$ with $3/2\leq s<2$ and $c\geq s+1$. The outsider has no profitable deviation: any quote that attracts positive-measure demand must fall below $s+1\leq c$, while a zero-demand quote yields zero profit. Consider a deviation $p$ by firm 1. Negative prices cannot improve profit. If $0\leq p\leq s-1$, total demand is at most three, so profit is at most $3(s-1)\leq3s/2$, the candidate profit. If $s-1<p\leq s$, the outsider remains foreclosed and \eqref{eq:duodemand} applies, giving
\[
 \pi_1(p)=p\left[\frac32+\frac34(s-p)\right].
\]
Its derivative is $3/2+3s/4-3p/2$, which is strictly positive throughout this interval because $s<2$. Hence no downward deviation is profitable. If $p>s$, the outsider becomes relevant on the long side and direct delivered-price comparison gives
\[
 q_1(p)=\max\left\{s+\frac32-p,0\right\}.
\]
On the positive-demand part, profit is $p(s+3/2-p)$. Its unconstrained maximizer $(s+3/2)/2$ is weakly below $s$ because $s\geq3/2$, so profit cannot increase once $p$ crosses $s$. Thus U1 is sufficient.

For U2, take $(2,2,r)$ with $r\geq3$ and $c\geq3$. The outsider cannot earn positive profit by attracting demand because positive demand requires a quote below three, weakly below its marginal cost. For a member deviation, removing the outsider can only weakly increase the deviator's demand and therefore gives an upper envelope for profit. Against the other member's price $2$, the two-member problem is globally maximized at $p=2$: for low prices profit is bounded by total market demand, on the interior two-member branch profit is the concave expression $p(3-3p/4)$ with unique maximizer $2$, and sufficiently high prices yield no gain. Hence U2 is sufficient.

Now let $(s,s,r)$ be any symmetric foreclosed Nash equilibrium of the unrestricted game. Foreclosure requires $r\geq s+1$, and outsider optimality requires $c\geq s+1$. A member price $s>2$ is impossible because a sufficiently small downward deviation remains on the regular foreclosed branch and raises profit. If $s<2$ and $r>s+1$, a sufficiently small upward member deviation also remains foreclosed and raises profit, so $r=s+1$ is necessary whenever $s<2$. Given $r=s+1$, a price $s<3/2$ is impossible because the right derivative after the outsider becomes relevant is $3/2-s>0$. Thus $3/2\leq s<2$ implies U1. If $s=2$, foreclosure and outsider optimality require $r\geq3$ and $c\geq3$, which is U2. These conditions exhaust the stated class and prove part 2.

For part 3, first note why the cost-floor restriction does not create additional equilibria within the stated class. At any cost-floor-feasible symmetric foreclosed candidate, each firm's candidate profit is nonnegative: the active member firms have zero marginal cost and nonnegative prices, while the foreclosed outsider earns zero. Any deviation deleted by the cost floor has $p_i<mc_i$; since demand is nonnegative, it yields $(p_i-mc_i)q_i\leq0$ and therefore cannot strictly improve on the candidate profit. Thus every symmetric foreclosed pure-strategy equilibrium of the cost-floor game is also a Nash equilibrium of the unrestricted game, so its candidate profiles are obtained by intersecting the U1--U2 set with the cost-floor restrictions.

The cost-floor game imposes $p_1,p_2\geq0$ and $p_3\geq c$. Intersect U1 with $r\geq c$. Since U1 also requires $r=s+1$ and $c\geq s+1$, both inequalities bind: $r=c=s+1$, so $s=c-1$. The conditions $3/2\leq s<2$ give $5/2\leq c<3$. Intersect U2 with $r\geq c$: the member price remains $2$, while the outsider may choose any $r\geq c$ for $c\geq3$. Restricting to $5/2<c<5$ yields \eqref{eq:piecewise-price}.
\end{proof}

The cost floor is an explicit modification of the strategy set, not a condition stated in \citet{gandalshy2001}, not an implication of Nash equilibrium, and not a claim that below-cost prices can be removed without loss of generality. Its role is narrower: within the symmetric foreclosed pure-strategy class, it removes the below-cost zero-sales outsider quotes that support the lower-price continuum and thereby pins down the member price in \eqref{eq:piecewise-price}.

For $5/2<c<3$ in the cost-floor game, the outsider has zero sales at $p_3=c$ but constrains the members' price at $c-1$: a higher member price restores outsider demand. In this limited static sense, the branch can be described as limit pricing or exclusion-maintenance pricing. This is not the informational signaling mechanism of \citet{milgromroberts1982}. At $c=3$ the two member-price branches meet at $2$ and the cost-based exclusion condition holds with equality; strict slack begins only for $c>3$.
